\scanpage{247}
\mspage{197}

et $=$\unsure{.} à priori des deux façons possibles

\diagram{\begin{tikzcd}[ampersand replacement=\&,row sep=1.4em]
\&C(C'(X))\&\\
C(X)\arrow[ur,"{C(i'_X)}"]\&\&C'(X)\arrow[ul,"{i_{C'(X)}}"']\\
\&X\arrow[ul,"i_X"]\arrow[ur,"i'_X"]\arrow[uu,dashed]\&,
\end{tikzcd}}

soit par $C(i'_X)\circ i_X$, soit par $i_{C(X)}\circ i'_X$, heureusement elles sont égales, car pour toute flèche $X\xrightarrow{i'}X'$ dans $\mathrm{Cat}_{/S}$, le carré

\diagram{\begin{tikzcd}[ampersand replacement=\&,row sep=.8em]
\&C(X')\&\\
C(X)\arrow[ur,"C(i')"]\arrow[dr,"i_X"']\&\&X'\arrow[ul,"i_{X'}"']\\
\&X\arrow[ur,"i'"]\&
\end{tikzcd}}

est commutatif (par fonctorialité de $i_X:X\mathbin{\unsure{-|\to}}C(X)$), et on applique ceci au cas de $X'=C'(X)$, $i'=i'_X$. \deleted{\unsure{Ce}}

Cette remarque nous encourage à faire un saut dans l'abstraction, et d'introduire la catégorie
% archive sheet number: 247
